\begin{definition}[Finite quotients of lattices and units]\label{mod:localfield-finitequotients}
Let $F$ be a nonarchimedean local field and write
$q_F=|k_F|$.  For every $m\leq n$ in $\mathbf Z$, the certified lattice
quotient $\mathfrak p_F^m/\mathfrak p_F^n$ has a
\texttt{Finite} instance and a specified noncomputable
\texttt{Fintype.ofFinite} enumeration, and
\[
 \left|\mathfrak p_F^m/\mathfrak p_F^n\right|
   =q_F^{(n-m).\operatorname{toNat}}.
\]
Both the \texttt{Nat.card} and the specified \texttt{Fintype.card} forms
of this equality are proved, as is the nonvanishing of the cardinality.
Likewise, for $m\leq n$ in $\mathbf N$,
$U_F^m/U_F^n$ has a \texttt{Finite} instance and a specified
noncomputable enumeration, with the exact boundary-sensitive formula
\[
 |U_F^m/U_F^n|=
 \begin{cases}
  1,&m=n=0,\\
  (q_F-1)q_F^{n-1},&m=0<n,\\
  q_F^{n-m},&0<m\leq n.
 \end{cases}
\]
The positive-depth formula, both depth-zero cases, the combined
\texttt{Nat.card} formula, and its specified \texttt{Fintype.card} form
are separate declarations; the unit-quotient cardinality is also proved
nonzero.

For $0<m\leq n$, the maps induced without selecting representatives by
\[
  [u]\longmapsto[u-1],\qquad [x]\longmapsto[1+x]
\]
are inverse equivalences
\[
 U_F^m/U_F^n\simeq \mathfrak p_F^m/\mathfrak p_F^n.
\]
The declarations constructing $u-1$ and $1+x$ record their exact
coercions to $F$, and the two quotient maps and the equivalence record
their values on quotient constructors.  This equivalence commutes with
denominator projection.  Under the additional linear-range
hypothesis $n\leq2m$, it is upgraded to a multiplicative equivalence
from $U_F^m/U_F^n$ to the multiplicative copy of the additive lattice
quotient; thus multiplication of unit classes is precisely addition of
their displacement classes.

The canonical quotient constructors compute addition, negation, and
subtraction on lattice representatives, and multiplication, inversion,
and division on unit representatives.  Projecting to the same denominator
is the identity, and successive denominator projections compose, for both
lattice and unit quotients.

Given a function on a numerator together with a proof that it is invariant
under depth-$n$ congruence, the declarations
\texttt{latticeQuotientLift} and \texttt{unitFiltrationQuotientLift}
descend it to the quotient, and their constructor equations recover the
original function.  The declaration
\texttt{latticeQuotientLift2} does the same for a two-variable function
separately invariant modulo its two denominators and has the corresponding
two-constructor equation.  Using the specified finite enumerations,
\texttt{latticeQuotientSum}, \texttt{latticeQuotientSum2}, and
\texttt{unitFiltrationQuotientSum} sum exactly these descended functions.
Translation by a lattice class and multiplication by a unit class are
proved to reindex the respective sums without changing their values.

If $m\leq n_1\leq n_2$, the kernel of
\[
 \mathfrak p_F^m/\mathfrak p_F^{n_2}\longrightarrow
 \mathfrak p_F^m/\mathfrak p_F^{n_1}
\]
is canonically equivalent to
$\mathfrak p_F^{n_1}/\mathfrak p_F^{n_2}$.  Every fiber is proved to have
$q_F^{(n_2-n_1).\operatorname{toNat}}$ elements.  The analogous unit
projection has kernel canonically equivalent to $U_F^{n_1}/U_F^{n_2}$;
every fiber has the cardinality given by exactly the
three-case unit-cardinality formula above with $(m,n)=(n_1,n_2)$.

Finally, suppose $m_i\leq n_i$ and
\[
 m_1+m_2\leq k,\qquad
 k\leq n_1+m_2,\qquad k\leq m_1+n_2.
\]
There is a representative-independent quotient product
\[
 (\mathfrak p_F^{m_1}/\mathfrak p_F^{n_1})\times
 (\mathfrak p_F^{m_2}/\mathfrak p_F^{n_2})
 \longrightarrow
 \mathfrak p_F^{m_1+m_2}/\mathfrak p_F^k,
 \qquad ([x],[y])\longmapsto[xy].
\]
The two upper bounds on $k$ ensure independence under changing either
representative, and the constructor equation is proved.  The product is
also bundled as an additive homomorphism in each variable, with an
evaluation theorem identifying it with the displayed product.  If
$\operatorname{ord}_F(a)=r$, division by $a$ canonically sends
$\mathfrak p_F^m/\mathfrak p_F^n$ to
$\mathfrak p_F^{m-r}/\mathfrak p_F^{n-r}$ and has constructor equation
$[x]\mapsto[x/a]$; it is also bundled as an additive homomorphism with its
evaluation theorem.  Product formation commutes with projection of either
input denominator and of its target denominator, and scalar division
commutes with denominator projection.  These are the quotient-arithmetic interfaces used to form
$[xy/\Gamma]$ in the finite-duality and stationary-pairing constructions.

\LeanFile{LanglandsFirstMainLemma/LocalField/FiniteQuotients.lean}
\lean{LanglandsFirstMainLemma.finiteLatticeQuotient}
\lean{LanglandsFirstMainLemma.instFiniteLatticeQuotient}
\lean{LanglandsFirstMainLemma.latticeQuotientFintype}
\lean{LanglandsFirstMainLemma.latticeQuotient_card}
\lean{LanglandsFirstMainLemma.latticeQuotient_fintype_card}
\lean{LanglandsFirstMainLemma.latticeQuotient_card_ne_zero}
\lean{LanglandsFirstMainLemma.positiveUnitDisplacement}
\lean{LanglandsFirstMainLemma.coe_positiveUnitDisplacement}
\lean{LanglandsFirstMainLemma.positiveUnitOfLattice}
\lean{LanglandsFirstMainLemma.coe_positiveUnitOfLattice}
\lean{LanglandsFirstMainLemma.positiveUnitFiltrationQuotientToLattice}
\lean{LanglandsFirstMainLemma.positiveUnitFiltrationQuotientToLattice_mk}
\lean{LanglandsFirstMainLemma.latticeQuotientToPositiveUnitFiltration}
\lean{LanglandsFirstMainLemma.latticeQuotientToPositiveUnitFiltration_mk}
\lean{LanglandsFirstMainLemma.positiveUnitFiltrationQuotientEquivLattice}
\lean{LanglandsFirstMainLemma.positiveUnitFiltrationQuotientEquivLattice_mk}
\lean{LanglandsFirstMainLemma.positiveUnitFiltrationQuotientMulEquivLattice}
\lean{LanglandsFirstMainLemma.positiveUnitFiltrationQuotientEquivLattice_projection}
\lean{LanglandsFirstMainLemma.finiteUnitFiltrationQuotient}
\lean{LanglandsFirstMainLemma.instFiniteUnitFiltrationQuotient}
\lean{LanglandsFirstMainLemma.unitFiltrationQuotientFintype}
\lean{LanglandsFirstMainLemma.positiveUnitFiltrationQuotient_card}
\lean{LanglandsFirstMainLemma.unitFiltrationQuotient_card_zero_zero}
\lean{LanglandsFirstMainLemma.unitFiltrationQuotient_card_zero}
\lean{LanglandsFirstMainLemma.unitFiltrationQuotient_card}
\lean{LanglandsFirstMainLemma.unitFiltrationQuotient_fintype_card}
\lean{LanglandsFirstMainLemma.unitFiltrationQuotient_card_ne_zero}
\lean{LanglandsFirstMainLemma.latticeQuotientMk_add}
\lean{LanglandsFirstMainLemma.latticeQuotientMk_neg}
\lean{LanglandsFirstMainLemma.latticeQuotientMk_sub}
\lean{LanglandsFirstMainLemma.unitFiltrationQuotientMk_mul}
\lean{LanglandsFirstMainLemma.unitFiltrationQuotientMk_inv}
\lean{LanglandsFirstMainLemma.unitFiltrationQuotientMk_div}
\lean{LanglandsFirstMainLemma.latticeQuotientProjection_self}
\lean{LanglandsFirstMainLemma.latticeQuotientProjection_trans}
\lean{LanglandsFirstMainLemma.unitFiltrationQuotientProjection_self}
\lean{LanglandsFirstMainLemma.unitFiltrationQuotientProjection_trans}
\lean{LanglandsFirstMainLemma.latticeQuotientLift}
\lean{LanglandsFirstMainLemma.latticeQuotientLift_mk}
\lean{LanglandsFirstMainLemma.latticeQuotientLift2}
\lean{LanglandsFirstMainLemma.latticeQuotientLift2_mk_mk}
\lean{LanglandsFirstMainLemma.unitFiltrationQuotientLift}
\lean{LanglandsFirstMainLemma.unitFiltrationQuotientLift_mk}
\lean{LanglandsFirstMainLemma.latticeQuotientSum}
\lean{LanglandsFirstMainLemma.latticeQuotientSum2}
\lean{LanglandsFirstMainLemma.unitFiltrationQuotientSum}
\lean{LanglandsFirstMainLemma.sum_latticeQuotient_add_left}
\lean{LanglandsFirstMainLemma.sum_unitFiltrationQuotient_mul_left}
\lean{LanglandsFirstMainLemma.latticeQuotientKernelEquiv}
\lean{LanglandsFirstMainLemma.latticeQuotientProjection_fiber_card}
\lean{LanglandsFirstMainLemma.unitFiltrationQuotientKernelEquiv}
\lean{LanglandsFirstMainLemma.unitFiltrationQuotientProjection_fiber_card}
\lean{LanglandsFirstMainLemma.latticeQuotientMul}
\lean{LanglandsFirstMainLemma.latticeQuotientMul_mk_mk}
\lean{LanglandsFirstMainLemma.latticeQuotientMulAddHom}
\lean{LanglandsFirstMainLemma.latticeQuotientMulAddHom_apply}
\lean{LanglandsFirstMainLemma.latticeQuotientDiv}
\lean{LanglandsFirstMainLemma.latticeQuotientDiv_mk}
\lean{LanglandsFirstMainLemma.latticeQuotientDivAddHom}
\lean{LanglandsFirstMainLemma.latticeQuotientDivAddHom_apply}
\lean{LanglandsFirstMainLemma.latticeQuotientProjection_mul}
\lean{LanglandsFirstMainLemma.latticeQuotientMul_projection_left}
\lean{LanglandsFirstMainLemma.latticeQuotientMul_projection_right}
\lean{LanglandsFirstMainLemma.latticeQuotientDiv_projection}
\leanok
\SourceText{Definition of Delta and Lemma lem:finite-duality.}
\uses{mod:localfield-unitfiltration,mod:localfield-residuefield}
\FormalizationContract No declaration selects a representative from a quotient.
Representative-level inputs only construct quotient classes, construct the positive-depth
displacement equivalence, or specify functions together with complete invariance proofs.
Character
perfectness, orthogonality, and the character-specific Delta and Lamprecht sums remain in their
downstream nodes.
\end{definition}
